% concatenated from main.tex
\documentclass[12pt,letterpaper,reqno]{amsart}
\usepackage{amsmath}
\usepackage{amsfonts}
\usepackage{amssymb}
\usepackage{graphicx}
\usepackage{xcolor}
\usepackage[dvipsnames]{xcolor}
\usepackage[english]{babel}
\usepackage{amsthm}

\headheight=8pt
\topmargin=0pt
\textheight=624pt
\textwidth=432pt
\oddsidemargin=18pt
\evensidemargin=18pt

\newenvironment{Proof}{\par\noindent{\it Proof.}} {\hfill$\scriptstyle\blacksquare$}
	
	
\subjclass[2020]{Primary 52A20, 53A07, 52B11.}
	
\keywords{Buoyancy, centroid, floating body, hedgehog, polytope.}


\allowdisplaybreaks
\newtheorem{theorem}{Theorem}
\newtheorem{lemma}{Lemma}
\newtheorem{claim}{Claim}
\newtheorem{remark}{Remark}
\newtheorem{definition}{Definition}
\usepackage{amsmath,amssymb}
\usepackage{amsfonts}
\usepackage{graphicx}
\usepackage[colorlinks=true, allcolors=blue]{hyperref}

%\def\subv{\delta |K|_d}           % volume submerged in the liquid
%\DeclareMathOperator{\vol}{vol} % volume
% \def\sof{P_{[\delta]}}    % Surface of flotation      \mathcal{C}_\delta K
% \def\bs{C_{\delta}}       % Surface of flotation      K_{[\delta]}
%$h_{K_{[\delta]}}(u)$      % ###### should be defined
%$h_{K}(u,\delta)$          % ###### should be defined

\begin{document}

\author{S. Dann}
\thanks{The first author is supported by the FAPA funds from Vicerrectoría de Investigaciones de la Universidad de los Andes (INV-2019-63-1699).}

\author{O. Herscovici}

\author{S. Myroshnychenko}
\thanks{The third author is supported in part by NSERC RGPIN-2024-05044.}

\address{S. Dann, Departamento de Matemáticas, Universidad de los Andes, Carrera 1 \#18A -- 12
Bogota, 111711, Colombia}
\email{s.dann@uniandes.edu.co}

\address{O. Herscovici, Department of Mathematics and Computer Science, St. John's University, 8000
Utopia Parkway, Queens, NY 11439, USA}
\email{herscovo@stjohns.edu}

\address{S.~Myroshnychenko, Department of Mathematics and Statistics,  University of the Fraser Valley, 33844 King Rd, Abbotsford, BC V2S 7M7,  Canada}
\email{serhii.myroshnychenko@ufv.ca}



\title[Uniqueness of Flotation and Buoyancy Surfaces for Convex Polytopes]{Uniqueness of Flotation and Buoyancy Surfaces for Convex Polytopes}
\maketitle



\begin{abstract}
We prove that a convex polytope \(P \subset \mathbb{R}^d\), \(d \ge 2\), of uniform density \(\delta \in (0,1)\) floating in a liquid of density \(1\), is uniquely determined by its surface of flotation \(P_{[\delta]}\) whenever \(\delta \neq \tfrac{1}{2}\). Analogously, we show that the buoyancy surface \(\mathcal{C}_\delta P\) of a convex polytope \(P\) with prescribed density \(\delta \in (0,1)\) uniquely determines~\(P\).
\end{abstract}



\section{Introduction}

Many questions of unique determination of convex bodies given the information about their sections or projections lie at the intersection of classical mechanics, harmonic analysis, differential geometry and classical convexity \cite{ARSY1,ARSY2,BMO,HSW,O,R23}. In this work,  we address the uniqueness questions of a convex polytope submerged in liquid based on information about its surface of flotation and the surface of buoyancy. 
Interest in these problems increased further after the striking discovery in \cite{R22} of a non-spherical body that floats in equilibrium in every orientation, thereby providing a negative answer to Ulam's question. Our results show that, within the polyhedral class, flotation and buoyancy surfaces encode remarkably rigid geometric information.


To describe our results, we briefly review the settings. Let \(K \subset \mathbb{R}^d\), \(d \geq 2\), be a convex body of uniform density \(\delta \in (0,1)\) floating in a liquid of uniform density \(\rho\). The condition \(\delta \in (0,1)\) is motivated by Archimedes' principle \cite{Ar}, which states that the submerged
volume \(V_{\mathrm{sub}}\) and the total volume \(|K|_d\) satisfy 
\[
\rho V_{\mathrm{sub}} = \delta |K|_d.
\]
After normalization, we may assume that \(\rho = 1\).

For each unit vector \(\theta \in \mathbb{S}^{d-1}\), there exists a unique
real number \(t_K(\theta,\delta)\) such that the liquid surface
$
H(\theta,t_K(\theta,\delta))
=
\{x \in \mathbb{R}^d : x \cdot \theta = t_K(\theta,\delta)\}
$
determines the submerged half-space
\[
H^-(\theta,t_K(\theta,\delta))
=
\{x \in \mathbb{R}^d : x \cdot \theta \ge t_K(\theta,\delta)\},
\]
for which
\[
|K \cap H^-(\theta,t_K(\theta,\delta))|_d
=
\delta |K|_d.
\]

\begin{figure}[h!]
    \centering
    \includegraphics[scale=0.15]{FS.png}
    \caption{Surface of flotation $K_{[\delta]}$ and buoyancy surface $\mathcal{C}_\delta K$ for a convex body $K$: $g(K \cap H(\theta))$ is the centroid of $K \cap H(\theta)$ and $g(\theta)$ is the centroid of $K \cap H^-(\theta)$.}
    \label{fig:FS}
\end{figure}

For the submerged cap
$
K \cap H^-(\theta,t_K(\theta,\delta)),
$
let \(g(\theta)\) denote its centroid. As \(\theta\) varies over the
unit sphere, the image of the map \(g\) traces a hypersurface
\(\mathcal{C}_\delta K\) in \(\mathbb{R}^d\), called the
\emph{surface of centers} (or \emph{buoyancy surface}) of \(K\),
see Figure~\ref{fig:FS}.

The envelope \(K_{[\delta]}\) of the truncating hyperplanes determined
by \(t_K(\theta,\delta)\) is the classical \emph{surface of flotation}
of \(K\). This surface has been extensively studied in affine convex
geometry and is closely related to affine surface area, affine
isoperimetric inequalities, and convex floating bodies
\cite{Du,SW}. By the Second Theorem of Dupin \cite{R23,Z}, the
contact point between \(K_{[\delta]}\) and the hyperplane
\(H(\theta)\) is the centroid \(g(K \cap H(\theta))\) of the section
\(K \cap H(\theta)\). We also recall from \cite{MR} that if \(K\) is
centrally symmetric, then \(K_{[\delta]}\) is the boundary of a convex
set and therefore coincides with the boundary of the convex floating
body \cite{W}.

Our first result establishes the injectivity of the map
\(\mathcal{C}_\delta\) within the polyhedral class of convex bodies.

\begin{theorem}\label{buoyTHM}
    Let $P$ and $Q$ be convex polytopes in $\mathbb{R}^{d}$, $d \geq 2$, with uniform densities $\delta_1, \delta_2 \in (0,1)$, respectively. 
    \begin{enumerate}
        \item If $d=2$, and  $\mathcal{C}_{\delta_1} P = \mathcal{C}_{\delta_2} Q,$ then $P= Q$ and $\delta_1 = \delta_2$.
        \item If $d \geq 3$, $\delta_1 = \delta_2 = \delta$, and $\mathcal{C}_{\delta} P = \mathcal{C}_{\delta} Q$, then $P=Q$.
    \end{enumerate}
    
\end{theorem}
Recall that if \(K\) is centrally symmetric, then
\(K_{[\frac12]}\) reduces to a single point, that is, the center of
symmetry of \(K\). Consequently, one cannot expect uniqueness from
the flotation surface at density \(\delta=\frac12\). In contrast,
for all other densities, we prove that the surface of flotation is
injective within the polyhedral class, even for possibly distinct
densities.
\begin{theorem}\label{floatingTHM}
    Let $P$ and $Q$ be convex polytopes in $\mathbb{R}^{d}$, $d \geq 2$, with uniform densities $\delta_1, \delta_2 \in (0,1)$, respectively. If at least one of $\delta_1, \delta_2$ differs from $\frac{1}{2}$, and
    $
    P_{[\delta_1]} = Q_{[\delta_2]},
    $
    then $P = Q$ and $\delta_1 = \delta_2$.
\end{theorem}

We also note that by the Third Theorem of Dupin \cite{R23,Z} for \textit{convex bodies} $P, Q \subset \mathbb{R}^2$, the \textit{simultaneous} equalities $P_{[\delta]} = Q_{[\delta]}$ and $\mathcal{C}_{\delta}P = \mathcal{C}_{\delta}Q$ imply $P = Q$. As will become clearer in the subsequent arguments (Claim \ref{bow-tie}), the equality of the buoyancy surfaces entails the equality of the corresponding chords of liquid in the bodies. Since the midpoints of these chords are precisely the contact points with the common surface of flotation, their endpoints trace out the same boundary, and hence the bodies coincide.

The proof of Theorem \ref{buoyTHM} in $\mathbb{R}^2$ takes advantage of the known structure of the buoyancy surface \cite{Z} and employs the Minkowski uniqueness theorem for polygons \cite{S}. In higher dimensions, by the Third Dupin Theorem \cite{R23, VP}, we show that the assumptions of the theorem imply parallel surfaces of flotation (Lemma \ref{buoy}), which yields subset inclusion for two polytopes, and their consequent equality (Lemma~\ref{inclusion}).

In dimensions $d=2$, the proof of Theorem \ref{floatingTHM} is based on considerations of the piecewise-hyperbolic structure of the flotation curve \cite{VP,Z}. In higher dimensions, the argument resembles the considerations in \cite{MR2018} and makes use of the piecewise-analytic structure of the flotation surface. Then the non-analytic points of the support function for $P_{[\delta]}$ on the unit sphere uniquely determine the vertices of $P$.

These results contribute to the broader program of identifying nonlinear geometric transforms that exhibit rigidity phenomena. 
It raises further questions concerning stability, extensions beyond the polyhedral class,  connections with the affine surface area, valuations, and other intrinsic volumes \cite{S}.



\section{Preliminaries}
We work in the ambient Euclidean space $\mathbb{R}^d$, $d \geq 2$, with the usual orthonormal basis $\{e_1, \ldots, e_d\}$. Given $x=(x_1, \ldots,x_d)$ and $y=(y_1, \ldots, y_d)$ in $\mathbb{R}^d$, their scalar product is $x\cdot y = \sum_{j=1}^d x_jy_j$, and the Euclidean norm $\|x\| = \left(x \cdot x\right)^{\frac 1 2}$. The unit sphere in $\mathbb{R}^d$ is $\mathbb{S}^{d-1} = \{x \in \mathbb{R}^d: \, \|x\|=1\}.$ 
For a differentiable function $h$ on $\mathbb{S}^{d-1}$, $\nabla_S h$ stands for its spherical gradient, that is, the projection of the Euclidean gradient onto the tangent space of $S^{d-1}$.
For a subset $S \subset \mathbb{R}^d$ and $\lambda \in \mathbb{R}$, its homothetic copy is
$
\lambda S = \{\lambda x: \, x \in S\},
$
where $\lambda x = (\lambda x_1,\lambda x_2, \ldots, \lambda x_d)$. Let $\partial S$ denote the boundary of $S$, and let $\text{int}S$ denote its interior. We also use conv$S$ to denote the convex hull of $S$. For a subspace $H \subset \mathbb{R}^d$, $H^\perp$ stands for its orthogonal complement.

A body $K \subset \mathbb{R}^d$ is a compact subset of $\mathbb{R}^d$ with non-empty interior. A body $K$ is convex if for every pair of points $x,y \in K$, the closed segment 
$$
[xy] = \{\lambda x + (1-\lambda)y: \, 0 \leq \lambda \leq 1 \}
$$
is contained in $K$. We use the notation $(xy)$ and $[xy)$ to denote the corresponding open and half-closed segments, respectively. 
We are interested in a special family of convex bodies -- namely, convex polytopes, which can be realized as the convex hull of a finite set of points in $\mathbb{R}^d$.
The centroid of a body $K$ is 
\begin{equation} \label{centre}
    g(K) = \frac{1}{|K|_d} \int_K x \, dx,
\end{equation}
where $| \cdot |_d$ stands for the $d$-dimensional Lebesgue measure (volume).
We assume that a convex body $K$ has a uniform density $\delta \in(0,1)$, so it floats in a liquid of density $1$. Then the volume of $K$ under the surface of the liquid is equal to $\delta |K|_d$. 

Instead of rotating the body within the liquid, we fix $K$ and rotate the liquid surface $H(\theta)$. These liquid surfaces give rise to the $t_K(\theta, \delta)$ mentioned in the Introduction, and we adopt the notation $h_{K_{[\delta]}}(\theta)$ for the signed distance from the origin to each plane along the normal vector $\theta \in \mathbb{S}^{d-1}$. For each surface of the liquid, the center of buoyancy is defined as the centroid of the
submerged part of $K$. The surface of buoyancy $\mathcal{C}_\delta K$ is the geometric locus of all centers of buoyancy. The surface of flotation $K_{[\delta]}$ is the envelope of all the surfaces of liquid. If $K_{[\delta]}$ is the boundary of a convex body, then it is the well-known convex floating body $K_\delta$, \cite{W}.




\section{Uniqueness of Flotation Surfaces $P_{[\delta]}$}
To prove Theorem \ref{floatingTHM}, we first show that $P_{[\delta]}$ is a $C^1$-smooth pieceswise-analytic hedgehog \cite{MM}.
\begin{lemma}\label{hedgehog}
Let $P$ be a polytope in $\mathbb{R}^d$, $d \geq 2$, with uniform density $\delta \in (0,1)$. 
Then $h_{P_{[\delta]}}$ is continuously differentiable and piecewise analytic on $\mathbb{S}^{d-1}$. 
The non-analytic points form a finite union of closed $(d-2)$-dimensional subsets of $\mathbb{S}^{d-1}$.
\end{lemma}


\begin{proof}
For any convex body $P$ and the liquid level $H(\theta) = \{x \in \mathbb{R}^d: \, x \cdot \theta = h_{P_{[\delta]}}(\theta)\}$ with the normal vector $\theta$, the centroid  $g_P(\theta)$ of the section $P \cap H(\theta)$ (see \eqref{centre}) varies continuously with $\theta \in \mathbb{S}^{d-1}$. By the Second Theorem of Dupin \cite[p.~287]{VP}, 
$$g_P(\theta)=\nabla_S h_{P_{[\delta]}}(\theta) + h_{P_{[\delta]}}(\theta) \theta, \quad g_P(\theta) = H(\theta) \cap P_{[\delta]},$$ 
Thus, $\nabla_S h_{P_{[\delta]}}(\theta)$ is continuous, i.e., $h_{P_{[\delta]}}$ is $C^1$-smooth.


By the approach in \cite{Nu}, let $\partial P$ be triangulated. 
Label the vertices of a triangulated facet $\partial P_i$ by 
$\{v_i^j\}_{j=1}^{d}$ in positive orientation and its outer unit normal vector by $N_i$. 
Then the $d$-dimensional volume $|P|_d$ of the polytope $P$ can be evaluated as
\begin{equation}\label{vol}
|P|_d = 
\frac 1 d \sum_i (v_i^j \cdot N_i) |\partial P_i|_{d-1} \quad \text{for any} \quad v_i^j \in \partial P_i.
\end{equation}
Labeling the vertices of a facet $\partial P_i$ (which is a $(d-1)$-dimensional simplex) in the order of the positive orientation, its volume can be found using the determinant of $(d-1) \times (d-1)$ matrix,
$$
|\partial P_i|_{d-1} = \frac{1}{(d-1)!} \begin{vmatrix}
     v_i^{2} - v_i^1\\ 
     v_i^{3}- v_i^1\\
     \ldots\\
     v_i^{d}- v_i^1\\
\end{vmatrix},
$$
where, without loss of generality, we assume that $\partial P_i$ lies in a horizontal hyperplane~$e_d^\perp$, so each row is a vector in $\mathbb{R}^{d-1}$. 
Since the vectors $ v_i^j-v_i^1, j=2,3,\ldots,d$ are linearly independent, the outer unit normal vector $N_i = \frac{\widetilde{N}_i}{|\widetilde{N}_i|}$ to $\partial P_i$ can be determined using the generalized vector product 
$$
\widetilde{N}_i = 
\begin{vmatrix}
     e_1 & \ldots & e_{d}\\ 
     &v_i^{2}- v_i^1\\
     &\ldots\\
     &v_i^{d}- v_i^1\\
\end{vmatrix},
$$
where now the rows are vectors in $\mathbb{R}^d$ with the last zero coordinate.
We apply the above considerations in the submerged halfspace $H^-(\theta)=\{x \in \mathbb{R}^d: \, x \cdot \theta \geq h_{P_{[\delta]}}(\theta)\}$ to the polytope $P(\theta)=P \cap H^-(\theta)$ with triangulated surface $\partial P(\theta)$. Next, we partition the sphere $\mathbb{S}^{d-1}$ into two sets. The first set contains the $\theta$'s with the property that a supporting hyperplane of $P_{[\delta]}$ perpendicular to them intersects only the interiors of edges of $P$, while the second set contains the remaining $\theta$'s. A supporting hyperplane to $P_{[\delta]}$, whose normal vector $\theta$ belongs to the second set, passes through a vertex of $P$.
 
Let us describe these two sets precisely. 
We enumerate the edges of $P$ with an index. 
For each $\theta \in \mathbb{S}^{d-1}$ in the first set, let $\mathcal{S_\theta}$ denote a finite set of indices for edges intersected by the liquid's surface $H(\theta)$. 
On the other hand, each index set $\mathcal{S_\theta}$ determines an open subset of $Y_\mathcal{S_\theta} \subset \mathbb{S}^{d-1}$ of directions for which the corresponding supporting hyperplanes intersect the same edges. 
In this way, $\mathbb{S}^{d-1}$ is split into a finite union of open sets $\cup_{\mathcal{S_\theta}}Y_\mathcal{S_\theta}$ and the rest is $X$, the boundaries of closures of $Y_\mathcal{S_\theta}$'s. 
For a vertex $v_i^j$, let 
\begin{align*}
X_{ij}=\Big\{\theta \in \mathbb{S}^{d-1}: \quad
h_{P_{[\delta]}}(\theta)-v_i^j\cdot  \theta =
0\Big\},
\end{align*}
then
$$
X = \bigcup_{v_i^j} X_{ij}.
$$


Hence $X_{ij}$ is a level set for the smooth map,
$$
F_{ij}(\theta):=h_{P_{[\delta]}}(\theta)-v_i^j\cdot  \theta, \quad F_{ij}: \mathbb{S}^{d-1} \to \mathbb{R}.
$$
Indeed, for the supporting point $g_P(\theta) = g(P \cap H(\theta))$, $g_P(\theta) \in P_{[\delta]} \cap H(\theta)$, we have
\begin{align*}
    \nabla_S F_{ij}(\theta) &= \nabla_S h_{P_{[\delta]}}(\theta) - \nabla_S(v_i^j \cdot \theta) =
    \nabla_S h_{P_{[\delta]}}(\theta) - (v_i^j - (v_i^j \cdot \theta) \theta )\\
    &=(g_P(\theta) - h_{P_{[\delta]}}(\theta) \theta) + (v_i^j \cdot \theta) \theta - v_i^j = (g_P(\theta) - v_i^j) + (v_i^j \cdot \theta - h_{P_{[\delta]}}(\theta) ) \theta.
\end{align*}
Since $(g_P(\theta) - v_i^j) \perp \theta $, we have 
$$
\| \nabla_S F_{ij}(\theta) \|^2 = \|g_P(\theta)-v_i^j\|^2 + |v_i^j \cdot \theta - h_{P_{[\delta]}}(\theta)|^2 \geq \|g_P(\theta)-v_i^j\|^2>0,
$$
as $h_{P_{[\delta]}}$ does not intersect the boundary of $P$.
Hence, for example, by (\cite{Hi}, Proposition 3.2, p.~22) level set $X_{ij}$ is a $(d-2)$-dimensional subset of $\mathbb{S}^{d-1}$.
Thus, for $P(\theta)$ we have the following,
\begin{enumerate}
    \item $|P(\theta)|_d = \delta|P|_d$ for any $\theta \in \mathbb{S}^{d-1}$.
    \item $P(\theta) =$ conv$\{v_1, \ldots, v_s, u_1(\theta), \ldots, u_t(\theta)\}$, where $\{v_k\}_{k=1}^s$ are vertices of $P$ from the half-space $x \cdot \theta > h_{P_{[\delta]} }(\theta)$, and $\{u_j(\theta)\}_{j=1}^t$ are vertices of $P$ obtained by intersecting the supporting hyperplane of $P_{[\delta]}$, $x \cdot \theta = h_{P_{[\delta]}}(\theta)$, with edges of~$P$.
\end{enumerate}
To determine $u_l(\theta)$, we solve the following system for $\lambda \in \mathbb{R}$,
$$
\begin{cases}
    x \cdot \theta = h_{P_{[\delta]}}(\theta)\\
    x = \lambda (v_{j_2} - v_{j_1}) + v_{j_1},
\end{cases}
$$
for some vertices $v_{j_1}, v_{j_2}$ of $P$ that span an intersected edge. This implies
$$
u_l(\theta) = \frac{h_{P_{[\delta]}}(\theta) - v_{j_1} \cdot \theta }{(v_{j_2} - v_{j_1})\cdot \theta} (v_{j_2} - v_{j_1}) + v_{j_1}.
$$
Substituting vertices $u_k$'s and $u_j(\theta)$'s into \eqref{vol}, for a given set $\mathcal{S}_\theta$, we obtain a level set of an analytic function,
$$
|P(\theta)|_d = \delta |P|_d, \quad \theta \in Y_\mathcal{S_\theta}.
$$
By the Implicit Function Theorem for analytic functions \cite{KP}, the above can be solved explicitly for $h_{P_{[\delta]}}=h_{P_{[\delta]}}(\theta)$ in some neighborhood $U$ of $\theta \in Y_\mathcal{S_\theta}$. 
This small neighborhood $U$ has the following properties: int $U \neq \emptyset$, and $U$ does not intersect $X$.

Thus, analytically extending $h_{P_{[\delta]}}(\theta)$ (and uniquely, as we intersect the same set of edges) in each connected component of $\mathbb{S}^{d-1} \setminus X$ yields an analytic function on $Y_\mathcal{S_\theta}$ for each possible index set $\mathcal{S_\theta}$.
\end{proof}

\subsection{Proof of Theorem \ref{floatingTHM} for $d =2$}
In dimension two, the flotation surface admits a piecewise-hyperbolic
description, allowing a direct reconstruction argument. In higher
dimensions, we instead exploit the analytic structure of the support
function and its singular locus.
\begin{lemma}\label{P_Dupin}
For a polygon $P$, $P_{[\delta]}$ is a finite union of hyperbolic arcs. At a common point $p$ of two consecutive hyperbolic arcs we have one of the following:
\begin{description}
    \item[Case 1] Hyperbolic arcs have a common tangent line (smooth union). This tangent line to $P_{[\delta]}$ at $p$ intersects two non-parallel sides of~$P$.
    \item[Case 2] Hyperbolic arcs have distinct tangent lines (non-smooth union). The set of supporting lines to $P_{[\delta]}$ at $p$ intersects a pair of parallel sides of $P$. 
\end{description}
\end{lemma}
\begin{proof}
    Recall that for the hyperbola $y=\frac k x, k>0$, the right triangle in the first quadrant bounded by any of its tangent lines and the coordinate axes has a constant area (of $2k$ square units). 
    Consider the family of lines that intersect two non-parallel sides of $P$ and cut out volume $V$ in $P$. Extend these two sides toward their intersection, and let the finite area between the lines in the exterior of $P$ be $V_0$. By means of a shear transformation, let these two lines be orthogonal. Then this family of lines cuts out the same area $V_0+V$, which means their envelope is a hyperbola.
    Also, if the family of lines intersects two parallel sides of $P$, then they all pass through one point. The intersection point belongs to $P_{[\delta]}$.
    
    The fact that only the above two cases can occur follows from the First Theorem of Dupin (a supporting line to $P_{[\delta]}$ is parallel to a tangent line to $\mathcal{C}_\delta P$, (\cite{VP}, p. 287)), and the following considerations on $\mathcal{C}_\delta P$.
    At a common point of two consequent hyperbolic arcs on $\mathcal{C}_\delta P$ the tangent line is parallel to the supporting line at $p \in P_{[\delta]}$ (Case 1).
    If $\mathcal{C}_\delta P$ contains a parabolic arc, then the set of all its tangent lines corresponds to the set of supporting lines to $P_{[\delta]}$ at a single point $p$ (Case 2).  
\end{proof}
The following lemma highlights the special settings for the case $\delta = \frac 1 2$.
\begin{lemma}\label{Sus}
    In $\mathbb{R}^d$, $d \geq 2$, $\delta \neq \frac{1}{2}$ if and only if for any point $x \in \partial K$ and $(d-2)$-dimensional subspace $\ell$, $x \in \ell$, $\ell \cap \text{int}K = \emptyset$, there exist two distinct hyperplanes that support $K_{[\delta]}$ and contain $\ell$.
\end{lemma} 

\begin{proof}
    Let $\delta \in \left(0, \frac 1 2 \right]$ and the two hyperplanes supporting $K_{[\delta]}$ split $K$ into three sets $K = K_1 \cup K_2 \cup K_3$ with $K_2$ enclosed by two of them. First, suppose $|K_1|_d=|K_3|_d=\delta |K|_d \leq \frac{|K|}{2}$. Then $|K_2|_d = |K|_d-2\delta|K|_d$. Clearly, the chords coincide, $|K_2|_d=0$, if and only if $\delta = \frac{1}{2}$. Now, if $\delta \in \left(\frac 1 2, 1 \right)$, we can repeat the same considerations for the complements of the above sets $K_1$ and $K_3$.

    Now, suppose for $\delta \in (0,1)$,  $$|K_2|_d+|K_3|_d = \delta |K|_d = |K_2|_d+|K_1|_d,$$ which yields $|K_1|_d = |K_3|_d = \delta|K|_d - |K_2|_d = |K|_d-\delta|K|_d$. Thus, $|K_2|_d = 2\delta|K|_d - |K|_d$, which equals zero if and only if $\delta = \frac{1}{2}$.
\end{proof}

\begin{remark}
We also note that the condition $\delta = \frac 1 2$ implies $h_{P_{[\delta]}}(-\theta) = -h_{P_{[\delta]}}(\theta)$ for all $\theta \in \mathbb{S}^{d-1}$. And if there exists at least one $\theta \in \mathbb{S}^{d-1}$ such that $h_{P_{[\delta]}}(-\theta) = -h_{P_{[\delta]}}(\theta)$, then $\delta = \frac 1 2$.
\end{remark}


The next lemma provides a billiard-like algorithm for the unique reconstruction of $P$ from $P_{[\delta]}$ in $\mathbb{R}^2$ and finishes the proof of Theorem \ref{floatingTHM} in $\mathbb{R}^2$.
\begin{lemma}
    For $\delta \neq \frac{1}{2}$ and a given $P_{[\delta]}$, there exists a unique polygon $P$ that can be reconstructed from~$P_{[\delta]}$.
\end{lemma}
\begin{proof}
    To reconstruct $P$, we first need to understand how the supporting lines of $P_{[\delta]}$ intersect sides of $P$. As Lemma \ref{P_Dupin} explains, $P_{[\delta]}$ is a finite union of hyperbolic arcs that are joined at the end points in a smooth or non-smooth way. We start by determining all asymptotes of the hyperbolic arcs of $P_{[\delta]}$, and considering the intersection $W$ of all quarter-planes containing the corresponding arcs of hyperbolas. Note that $W$ is also a polygon. By the above observation, $P \subseteq W$ and $\partial P \cap \partial W \neq \emptyset$, where the intersection consists of a finite number of segments. Note that by convexity, each subset of $\partial P \cap \partial W$ that belongs to one line can be assumed to be a segment.
    
    To determine missing sides of $P$, we plot the one-sided tangent rays to non-smooth adjoint hyperbolic arcs at the common points. The sides of $W$ intersected by these two rays must contain the end points of missing sides; the segments between the intersected points on such sides of $W$ are subsets of $\partial P$.
    
    Observe that the intersection of a supporting line of $P_{[\delta]}$ with $\partial P$ is continuous with respect to a direction vector $\theta \in \mathbb{S}^2$.
    We partition each side of $P$ into a finite number of open segments: we exclude all vertices and all points on the sides that lie on liquid levels passing through vertices.
    The number of such excluded points is bounded by three times the number of vertices of $P$ (since each vertex may exclude at most two distinct points).
    This divides $\partial P$ into a finite number of segments. 
    Thus, for any point $x$ in a given segment $(pq)$, two supporting lines of $P_{[\delta]}$ that contain $x$ intersect the same pair of segments (which may coincide) regardless of the choice of $x \in (pq)$.

    
    For $x \in (pq)$, by Lemma~\ref{Sus}, let $y \in (rs)$ and $z \in (uv)$ be the two distinct points where the lines of liquid through $x$ intersect $\partial P$. 
    Then we investigate the following cases based on the mutual position of the intersected segments $(pq), (rs), (uv)$.

       \begin{figure}[h!]
        \centering
        \includegraphics[scale=0.35]{parallel-cases.png}
         \caption{For $t \in P_{[\delta]}$, case $(rs)\neq (uv)$ and $pq \parallel (uv)$.}
    \label{fig:2DParallel}
     \end{figure}
     
    \begin{itemize}
     \item If $(rs)\neq (uv)$ and $(pq) \not \parallel (uv)$, then $(pq), (uv) \subset \partial W$.

     \item If $(rs)\neq (uv)$ and $(pq) \parallel (uv)$, then $(pq)$ may be parallel to $(rs)$, Figure~\ref{fig:2DParallel}. Then, $(rs)$ and $(uv)$ belong to the same side of $P$ (otherwise, they are not parallel and belong to $\partial W$). 
     If we let $x$ approach $q$ (then $y$ approaches $\tilde{y} \in [rs)$), then the second line of liquid through $y$ may intersect the same side to which $(pq)$ belongs.
     
     The intersection cannot be a point $w_1$ to the left from $x$ along the line through $p, q$, since then the area to the left from the line through $w_1, y$ is strictly less than $\delta$ (less than the area on the left from line through $x, y$). 
     If the intersection is at a point $w_2 \in (pq]$, then when $x$ approaches $q$, $w_2$ also approaches $q$, so $x=w_2=q$ in the limit. However, Lemma \ref{Sus} guarantees $x\neq w_2$. 
     Thus, the point is in a position of $w_3$ to the right from $q$ along the line through $p,q$.
    Consequently, $w_3$ lands in a segment $(p_1q_1) \neq (pq)$ on the right from $(pq)$. Hence, we can repeat the same considerations starting from $(p_1q_1)$ until we exhausted all segments laying on the same side as $(pq)$ or $(rs)$. 
    As the number of such segments on each side is finite, the point corresponding to $\tilde{y}$ or $w_3$ will belong to a segment not parallel to $(pq)$.
    Hence, that segment and its non-parallel predecessor $(s_1s_2)$ belong to $\partial W$. Based on $P_{[\delta]}$ and a given segment $(s_1s_2) \in \partial P$, we reconstruct the previous segments all the way to $(pq)$.
    

     \item If $(rs)=(uv)$ and $(pq) \not \parallel (uv)$, analogously $(pq), (uv) \subset \partial W$

     \item If $(rs)= (uv)$, and $(pq) \parallel (uv)$, we can apply the same approach as $(rs) \neq (uv)$ and $(pq) \parallel (uv)$, because the position of $z \in (rs)$ does not affect the argument for $x$ approaching $q$.
    \end{itemize}

We conclude that by continuously varying the liquid line to intersect $\partial W$, we reconstruct $\partial P$, except for a finite number of points that we excluded. Taking the closure of this set yields $\partial P$, and $P$.
\end{proof}

We also remark that even for non-centrally-symmetric polytopes density $\delta=\frac 1 2$ does not yield uniqueness.
\begin{remark} \label{EX}
\begin{figure}[h!]
    \centering
    \includegraphics[scale=0.2]{delta-nuance.png}
    \caption{$P_{[\frac 1 2]}$ does not define non-symmetric polytopes uniquely.}
    \label{delta_nuance}
\end{figure}
Let $P$ be a non-symmetric polygon of density $\delta= \frac 1 2$ that contains two parallel sides of equal length. Then construct a polygon $Q$ of the same density $\frac 1 2$ that consists of the same sides, except that the two parallel sides are replaced with four pairwise parallel sides of pairwise equal lengths and the corresponding enclosed areas coincide, $A_1=A_2$ (Figure~\ref{delta_nuance}). In this scenario, $P_{[\frac 1 2]} = Q_{[\frac 1 2]}$, yet $P \neq Q$.
\end{remark}

\subsection{Proof of Theorem \ref{floatingTHM} for $\mathbb{R}^d, d \geq 3$}
Now the proof of Theorem 2 relies on the observation that the support
function of the flotation surface is analytic except in directions
corresponding to liquid levels passing through vertices of
the polytope (Lemma \ref{hedgehog}). Thus, the singular set of the support function encodes
the vertex structure of the polytope.


\begin{proof}
By Lemma \ref{hedgehog}, the support function $h_{P_{[\delta]}}$ is analytic on $\mathbb{S}^{d-1}$  except at points $\theta$ from the following set,
$$
Z=\left\{\theta \in \mathbb{S}^{d-1}: \quad  x \cdot \theta = h_{P_{[\delta]}}(\theta) \quad \text{contains a vertex of} \, P\right\}.
$$
Thus, every vertex of $v \in P$ is in bijective correspondence with $Z_v$,
$$
Z_v=\left\{\theta \in \mathbb{S}^{d-1}: \quad  v \cdot \theta = h_{P_{[\delta]}}(\theta)\right\}.
$$

If there was another polytope $Q \neq P$ such that the condition of Theorem \ref{floatingTHM} is satisfied, we would have a vertex $w \in Q$ but not a vertex of $P$, such that
$$
Z_w \in Z, \quad w \text{\, is not a vertex of \,} P,
$$
which is impossible. In other words, for $d \geq 3$, the non-analytic points of $h_{P_{[\delta]}}$ on $\mathbb{S}^{d-1}$ uniquely reconstruct $P$ by identifying the supporting cones of $P_{[\delta]}$ corresponding to the vertices of $P$.


A vertex $v \in P$ cannot be reconstructed in this way when the planes corresponding to $Z_v$ do not intersect (the cone becomes a cylinder). This is possible only when $Z_v$ is a large subsphere on $\mathbb{S}^{d-1}$. In this case, the planes intersect in an $(d-2)$-dimensional affine subspace $\ell$ passing through $v$ and int$P$. For any $\theta$ perpendicular to $\ell$, the liquid surfaces corresponding to $\theta$ and $-\theta$ coincide, resulting in $\delta = \frac 1 2$.
\end{proof}



\section{Uniqueness of Buoyancy Surfaces $\mathcal{C}_\delta P$}

\subsection{Proof of Theorem \ref{buoyTHM} in $\mathbb{R}^2$}

For a polygon $P$ in $\mathbb{R}^2$, $\mathcal{C}=\mathcal{C}_\delta P=\mathcal{C}_\delta Q$ is a closed convex curve that consists of hyperbolic and parabolic arcs (\cite{Z}, \cite{VP}). By construction, for an arc $a$ of $\mathcal{C}$ there exist a pair of intersected sides $p_1^a, p_2^a$ of $P$ and the respective sides $q_1^a, q_2^a$ of $Q$.

\begin{claim}
Side $p_1^a$ is parallel to $q_1^a$ or $q_2^a$.
\end{claim}

\begin{proof}
    If $a$ is a hyperbolic arc, sides $p_1^a, p_2^a\in P$ are known to be parallel to the asymptotes of the hyperbola that contains the arc $a$ (\cite{Z}, p. 670 ). Likewise for $q_1^a,q_2^a\in Q$. Thus, $p_1^a\parallel q_1^a$ or $p_1^a\parallel q_2^a$. The correspondence between the sides can be determined uniquely due to the given position of the hyperbola, it belongs to a quarter plane determined by its intersecting asymptotes.

    If $a$ is a parabolic arc, the sides $p_1^a, p_2^a\in P$ are parallel to the axis of symmetry of $a$, (\cite{Z}, p. 673 ) Likewise for $Q$. Then all four sides are parallel, $p_1^a\parallel p_2^a\parallel q_1^a\parallel q_2^a$. In this case, the correspondence between the sides is determined uniquely due to the given position of the parabola enclosed in the slabs determined by the pairs of the sides.
\end{proof}

\begin{claim} \label{parallel-sides}
Polygons $P$ and $Q$ have the same number of sides that are pairwise parallel.
\end{claim}

\begin{proof}
    Let $p_i, p_j$ be two sides of $n$-gon $P$ that give rise to an arc $a_{ij}$ of $\mathcal{C}$. Similarly, there is also a pair of sides $q_i,q_j$ of $Q$ corresponding to the same arc $a_{ij}$.
    Label the sides so that $p_i\parallel q_i$. Hence, to each side of $P$ we can associate a unique parallel side of $Q$. Moreover, adjacent sides are mapped onto adjacent sides. 
    Thus, $Q$ is also an $n$-gon. Therefore, polygons $P$ and $Q$ have the same number of sides, and  the sides with the corresponding indices are parallel, $q_i\parallel p_i$.
\end{proof}

\begin{claim}\label{bow-tie}
    For all $i=1,2, \ldots,n$, the length of $p_i$ is proportional by the same positive constant to the length of $q_i$.
\end{claim}

\begin{proof}
    For a hyperbolic arc in $\mathcal{C}$, there exists a pair of non-parallel sides in $P$ and the corresponding pair of non-parallel sides in $Q$. Denote this hyperbolic arc by $a$ and choose a point $A$ from the interior of $a$ with the tangent line $\ell(A)$. Then, by the First Dupin's Theorem (\cite{VP}, p. 287), this tangent line is parallel to liquid surfaces $\ell^P_1$ and $\ell^Q_1$ that cut out $\delta_1 |P|_2$ from $P$ and $\delta_2|Q|_2$ from $Q$.
    Likewise, for a point $B$ close to $A$ on $a$, Figure \ref{fig:2D_hyperbolic_arc}.
    
    Let $K_1, L_1$ be the points of intersection of $\ell^P_1$ with the sides of $P$, and $K_2,L_2$ be the points of intersection of $\ell^P_2$ with the sides of $P$. 
    For $K_1,K_2,L_1, L_2\in P$, we have the corresponding points $T_1, T_2, X_1, X_2 \in Q$, respectively. Therefore, we have $K_1L_1\parallel T_1X_1$, $K_2L_2\parallel T_2X_2$ and $K_1K_2\parallel T_1T_2$, $L_1L_2 \parallel X_1X_2$.
    
    \begin{figure}[h!]
        \centering        \includegraphics[scale=0.3]{Bow-ties}
        \includegraphics[scale=0.3]{Bow-ties-2}
        \caption{Hyperbolic arcs and two possible positions of the chords.}
        \label{fig:2D_hyperbolic_arc}
    \end{figure}

    Let $O_P$ be the point of intersection of lines $\ell_1^P$ and $\ell_2^P$, and $O_Q$ be the analogous point for lines $\ell_1^Q$ and $\ell_2^Q$. By the Third Theorem of Dupin (\cite{VP}, p. 288), for density $\delta$, the radius of curvature $R$ of $a$ is given by,
   \begin{equation}
       R=\frac{I}{\delta_1 |P|_d},
    \end{equation}
    where $I$ is the moment of inertia for the liquid level in the body (segment $K_1L_1$) around its centre of mass (midpoint of the segment $K_1L_1$).
    Choose a coordinate system such that $K_1L_1$ is parallel to the $x$-axis. Then the moment of inertia about the center of mass of $K_1L_1$ at the origin is
    \begin{equation}
        I=\int\limits_{-\|\frac{K_1L_1}{2}\|}^{\|\frac{K_1L_1}{2}\|} x^2\,dx=\frac{2x^3}{3}\,\Bigg{|}_{0}^{\|\frac{K_1L_1}{2}\|}=\frac{\|K_1L_1\|^3}{12}.
    \end{equation}
    Then we obtain,
    \begin{equation*}
        \frac{\| K_1L_1\|^3}{ \delta_1  |P|_2}=R =\frac{\| T_1X_1\|^3}{\delta_2 |Q|_2}.
    \end{equation*}
    Hence, 
    $$
    \| K_1L_1\|=\tilde{c} \| T_1X_1\|, \quad \|K_2L_2\|=\tilde{c} \| T_2X_2\|, \quad \text{for} \quad \tilde{c} = \sqrt[3]{\frac{\delta_1 |P|_2}{\delta_2  |Q|_2}}.
    $$
    Note that in a given quarter between a pair of intersecting lines, the length of parallel segments intersecting the lines is a monotonic function of the distance from the apex of the angle.
    
    In the hyperbolic case, the two pairs of lines containing the sides of $P$ and $Q$ are pairwise parallel, so congruence of the angles between these lines and $\tilde{c} \|T_1X_1\|=\|K_1L_1\|$ implies homothety of the sides, $\|L_1L_2\| = \tilde{c}  \|X_1X_2\|$ and, likewise, $\|K_1K_2\|=\tilde{c} \|T_1T_2\|.$

    For a parabolic arc, with the analogous notation, Figure \ref{fig:2D_parabolic_arc},
    $$
    K_1K_2 \parallel L_1L_2 \parallel X_1X_2 \parallel T_1T_2,
    $$
    and the areas of the triangles in each polygon are respectively equal,
    $$
    |\triangle K_1K_2O_P|_2 = |\triangle L_1L_2O_P|_2=\tilde{c}^2 |\triangle T_1T_2O_Q|_2, \quad |\triangle T_1T_2O_Q|_2 = |\triangle X_1X_2O_Q|_2.
    $$
    
    \begin{figure}[h!]
        \centering      
        \includegraphics[scale=0.3]{Bow-ties-3}
        \caption{Parabolic arcs and possible positions of the chords.}
        \label{fig:2D_parabolic_arc}
    \end{figure}
    
    Then, $O_P \in P_{[\delta_1]}$ has equal distance to both lines containing segments $K_1K_2$ and $L_1L_2$. Likewise for $O_Q$ and $X_1X_2$, $T_1T_2$. The equality of areas then implies $\|K_1K_2\| = \|L_1L_2\|$ and $\|T_1T_2\| = \|X_1X_2\|$. 
    Analogously to the hyperbolic case, $\|K_1L_1\| = \tilde{c} \|T_1X_1\|$. Along with the fact $O_p$ splits $L_1K_1$ and $K_2L_2$ in half, we also get $\|L_1L_2\| = \tilde{c} \|X_1X_2\|$, $\|K_1K_2\| = \tilde{c} \|T_1T_2\|$.
    \end{proof}
    
\begin{claim}\label{shift}
Polytopes coincide, $P = Q$.
\end{claim}

\begin{proof}
By Minkowski uniqueness theorem (\cite{S}, Section 8.2) for polygons, 
$$
P = \pm \tilde{c} \cdot Q + u, \quad u \in \mathbb{R}^2.
$$
Consider the two possible sign choices separately, denote $\tilde{Q} = \pm \tilde{c} Q$, and assume $\tilde{c} \neq 1$. Then $x_0 = \frac{1}{1\pm \tilde{c}}u$ is the fixed point of the affine map $T(x) = \pm \tilde{c}x+u$. 
If $P=\tilde{Q}+u$, then $\mathcal{C}=\tilde{c}\, \mathcal{C}+u$, and $\mathcal{C}-x_0 = \tilde{c} (\mathcal{C}-x_0)$. Thus, up to translation, $\mathcal{C} = \pm \tilde{c} \, \mathcal{C}$. Repeated application of the identity and the fact that $\mathcal{C}$ is not a point yield $\tilde{c} = 1$.

If $P = Q+u$, then $\mathcal{C} = \mathcal{C} + u$. Repeated application of this identity shows $u=0$.

If $P=-Q+u$, then $\mathcal{C}=-\mathcal{C}+u$, or equivalently $\mathcal{C} - \frac{1}{2}u = -\left(\mathcal{C} - \frac 1 2 u\right)$.
If we adjust the coordinate system to have the origin at $\frac 1 2 u$, then, by the above considerations, each side of $P$ has a parallel counerpart of the same length in $P$, which yields $-P=P=Q$.
\end{proof}

   
\subsection{Proof of Theorem \ref{buoyTHM} in $\mathbb{R}^d, d \geq 3$}

We begin with a lemma whose conclusion is likely not optimal, but sufficient for our purposes.
\begin{lemma}\label{buoy}
    Let $P$ and $Q$ be two convex bodies in $\mathbb{R}^{d}$, $d\geq 3$, with uniform densities $\delta_1, \delta_2 \in (0,1)$.  If their buoyancy surfaces coincide, 
    $$\mathcal{C}_{\delta_1} P = \mathcal{C}_{\delta_2} Q,$$
    then their surfaces of flotation are parallel surfaces,
    $$
    h_{Q_{[\delta_2]}}=h_{P_{[\delta_1]}} + c, \quad \text{for some} \quad c \in \mathbb{R}.
    $$
\end{lemma}

\begin{proof}
By the Third Theorem of Dupin (\cite{VP}, p. 288), the equality of buoyancy surfaces implies that for any $\xi \in \mathbb{S}^{d-1}$,
$$
\frac{1}{\delta_1|P|_d} I_{P \cap H_P(\xi)} (\Pi_P) = \frac{1}{\delta_2|Q|_d} I_{Q \cap H_Q(\xi)} (\Pi_Q),
$$
where
the surface of the liquid $H_P(\xi)$ is parallel to the tangent plane to $\mathcal{C}_\delta P$ with the normal vector $\xi$ (so it is parallel to $H_Q(\xi)$ as well);
$\Pi_P$ is any $(d-2)$-dimensional plane passing through the centroid $g_P(\xi)$ of $P\cap H_P(\xi)$ (and it is parallel to $\Pi_Q$); and
$$
I_{P \cap H_P(\xi)} (\Pi) = \int_{P\cap H_P(\xi)}\text{dist}^2(\Pi,v) \, dv.
$$
Here, dist denotes the Euclidean distance between a point $x$ and an affine subspace~$\Pi$. Since liquid levels are parallel, we may introduce the same coordinate system $\{x_1, \ldots,x_{d-1}\}$ for both $H_P(\xi)$ and $H_Q(\xi)$. Let the equation of $\Pi_P$ in these coordinates be
$$
\Pi_P: \, a_1x_1+\ldots+a_{n-1}x_{n-1} + c_P = 0, \qquad a:=(a_1,\ldots,a_{d-1}) \in \mathbb{S}^{d-2} \subset H_P(\xi).
$$
The fact that $g_P(\xi) \in \Pi_P$ implies $a \cdot g_P(\xi) = -c_P$.
We also note
$$
\text{dist}^2(\Pi_P,v) = \frac{1}{a\cdot a}(a \cdot v + c_P)^2 = (a\cdot v)^2 + 2(a \cdot v) c_P + c_P^2.
$$
Then,
\begin{align*}
&\frac{1}{\delta_1|P|_d} I_{P \cap H_P(\xi)} (\Pi_P) =\\
&\frac{1}{\delta_1|P|_d} \int_{P\cap H_P(\xi)}(a \cdot v)^2 \, dv + 2c_P |P \cap H_P(\xi)|_{d-1} \, (a \cdot g_P(\xi))  + c_P^2 |P \cap H_P(\xi)|_{d-1} = \\
&\frac{1}{\delta_1|P|_d} \cdot \sum_{i,j=1}^{d-1} a_i a_j \left(\int_{P\cap H_P(\xi)} v_i v_j \, dv\right) - \frac{c_P^2}{\delta_1|P|_d} |P \cap H_P(\xi)|_{d-1}. 
\end{align*}
Thus, for a fixed $\xi \in \mathbb{S}^{d-1}$,
$\frac{1}{\delta_1|P|_d} I_{P \cap H_P(\xi)} (\Pi_P) = \frac{1}{\delta_2|Q|_d} I_{Q \cap H_Q(\xi)} (\Pi_Q)$ holds for all $a \in \mathbb{S}^{d-2}$ (note that $\Pi_P$ is parallel to $\Pi_Q$ if and only if the corresponding $a$'s are parallel for both planes $\Pi_P$ and $\Pi_Q$). The equality of the quadratic forms above for any $a \in \mathbb{S}^{d-2}$ implies
$$
\frac{c_P^2}{\delta_1 |P|_d}  |P \cap H_P(\xi)|_{d-1} = \frac{c_Q^2}{\delta_2  |Q|_d}  |Q \cap H_Q(\xi)|_{d-1}.
$$
Note that $c_P$ is the distance from $\Pi_P$ to the centroid $P \cap H_P(\xi)$. However, if we choose the coordinate system such that $g_P(\xi) = (0,\ldots,0) \in \mathbb{R}^{d-1}$, then $c_P = 0$, which implies $c_Q = 0$, and $g_Q(\xi) = (0,\ldots,0) = g_P(\xi)$. Consequently, the segment connecting $g_P(\xi)$ and $g_Q(\xi)$ in $\mathbb{R}^{d}$ is always perpendicular to the liquid levels or the centroids coincide (if and only if the liquid surfaces coincide). If they do not coincide, we may regard $g_P(\xi)$ as the unique point the corresponding flotation surface and the liquid surface have in common. Thus, for a continuous function $d(\xi)$ and contact points $g_Q(\xi), g_P(\xi)$ of the corresponding envelopes, we have (see \cite{MM}),
\begin{align*}
    g_Q(\xi) &= g_P(\xi) + d(\xi) \xi\\
    h_{Q_{[\delta_2]}}(\xi) \xi  + \nabla_S h_{Q_{[\delta_2]}}(\xi) &= h_{P_{[\delta_1]}}(\xi) \xi  + \nabla_S h_{P_{[\delta_1]}}(\xi)+ d(\xi) \xi.
\end{align*}
Since $\nabla_S h_{P_{[\delta_1]}}(\xi) \perp \xi$ and $\nabla_S h_{Q_{[\delta_2]}}(\xi) \perp \xi$, we get 
$
\nabla_S h_{Q_{[\delta_2]}}(\xi) = \nabla_S h_{P_{[\delta_1]}}(\xi),
$
which implies $h_{Q_{[\delta_2]}}(\xi) = h_{P_{[\delta_1]}}(\xi) +c$ for any $\xi \in \mathbb{S}^{d-1}$ and a constant $c \in \mathbb{R}$. Consequently, $d(\xi) \equiv c$. 
\end{proof}

The following Lemma shows that two convex bodies $P \subseteq Q$ of same density $\delta$ coincide if and only if $\mathcal{C}_{\delta}P = \mathcal{C}_{\delta}Q$.
\begin{lemma}\label{inclusion}
Let $P$ and $Q$ be convex bodies in $\mathbb{R}^d, d \geq 3$, both of uniform density $\delta \in (0,1)$. If $P \subseteq Q$ and $\mathcal{C}_{\delta} P = \mathcal{C}_{\delta} Q$, then $P=Q$.
\end{lemma}

\begin{proof}

\begin{figure}[h]
    \centering
    \includegraphics[scale=0.22]{lem6.png}
    \caption{Functions $V_P, V_Q$ and inclusion $P \subseteq Q$.}
    \label{fig:lem6}
\end{figure}

For $h \in \mathbb{R}$, denote $P(\theta)=P \cap \{ x \cdot \theta \leq h\}$ and $V_P(h) = \left|P(\theta) \right|_d$. Since $P \subseteq Q$, we have $P(\theta) \subseteq Q(\theta)$, so $V_P(h) \leq V_Q(h)$ for $h \in \mathbb{R}$.
Also, $\delta|P|_d \leq \delta |Q|_d$, and $V_Q(h_{P_{[\delta]}}(\theta)) \geq V_P(h_{P_{[\delta]}}(\theta))$, so $V_Q(h) \geq \delta|P|_d$ for $h \geq h_{P_{[\delta]}}(\theta)$.
Then, $h_{P_{[\delta]}}(\theta) \leq h_{Q_{[\delta]}}(\theta)$, since $V_P(h_{P_{[\delta]}}(\theta)) \equiv~\delta |P|_d$, Figure \ref{fig:lem6}.

On the other hand, the centroid of $P(\theta) \setminus Q(\theta)$ is below the plane $x \cdot \theta = h_{Q_{[\delta]}}(\theta)$ in the direction $\theta$, while the centroid of $Q(\theta) \setminus P(\theta)$ is above the same plane in the direction $-\theta$. Also, since $\mathcal{C}_{\delta} P = \mathcal{C}_{\delta} Q$, we have $P(\theta) \cap Q(\theta) \neq \emptyset$. Thus, centroid of $P(\theta)$ should be strictly lower (along $\theta$) than the centroid of $P(\theta) \cap Q(\theta)$, but the opposite must hold for the centroid of $Q(\theta)$. This implies $P(\theta) = Q(\theta)$ for all $\theta \in \mathbb{S}^{d-1}$.
\end{proof}

\begin{figure}[h!]
    \centering
\includegraphics[scale=0.2]{Polytope-Inclusion.png}
    \caption{Inclusion of projections, $\tilde{P} \subseteq \tilde{Q}$.}
    \label{fig:incl}
\end{figure}

To conclude the proof of Theorem \ref{buoyTHM}, assume that $c \geq 0$ in Lemma \ref{buoy}. Then, by Lemma \ref{inclusion}, our goal is to show $P \subseteq Q$.  
Let $\ell$ be an $(d-2)$-dimensional plane that contains a vertex $u \in P$, $\ell \cap \text{int}P = \emptyset$, and $W_1, W_2$ be two hyperplanes containing $\ell$ that support $P_{[\delta]}$. 
The spanning vectors of $\ell$ for all vertices are dense on $\mathbb{S}^{d-1}$, and exclude the vectors parallel to the facets of the polytopes. 
Consider the projections $\tilde{u}$ of $u$, $\tilde{P}$ of $P$, and $\tilde{P}_{[\delta]}$ of $P_{[\delta]}$ onto $\ell^\perp$, dim$(\ell^\perp) = 2$, Figure~\ref{fig:incl}.

The projections of $W_1$ and $W_2$ onto $\ell^\perp$ are lines $w_1, w_2$ that support $\tilde{P}_{[\delta]}$. Since $\tilde{P}_{[\delta]}$ is connected, the ray from $\tilde{v}$ through $\tilde{u}$ intersects $\tilde{P}_{[\delta]} \subset \tilde{P}$ at a point $y$.
Indeed, for vertex $\tilde{v}$ of $\tilde{Q}$, 
$$
\tilde{v} = \tilde{u} +\frac{c}{\sin\alpha} \frac{\tilde{u}-y}{\|\tilde{u}-y\|}, \qquad \tilde{u}, y \in P,
$$
where $0<\alpha< \pi$ is the angle between $y-\tilde{u}$ and $w_2$.
Hence, $\tilde{P} \subseteq \tilde{Q}$ for all $\ell$ not parallel to the facets of the polytopes. Such set consists of almost all points on $\mathbb{S}^{d-1}$, thus $P \subseteq Q$.

\section*{Acknowledgment}
We are grateful to Dmitry Ryabogin for the inspiration and helpful discussions.



\begin{thebibliography}{9}
\bibitem{Ar}
Archimedes, The works of Archimedes, translated by T. Heath, Dover 2002.

%\bibitem{ANRY}
%M. Alfonseca, F. Nazarov, D. Ryabogin, V. Yaskin, ``\textit{A solution to the fifth and the eighth Busemann-Petty problems in a small neighborhood of the Euclidean ball}", Adv. Math., \textbf{390} (2021), Paper No. 107920, 28 pp.

\bibitem{ARSY1}
M. Alfonseca, A. Stancu, D. Ryabogin, V. Yaskin, ``\textit{On flotation, stability and related questions: A survey}", to appear in New Probes into Discrete and Convex Geometry (J. Pach and G. Toth, eds.), Bolyai Society Mathematical Studies, Vol. \textbf{32}, Springer-Verlag, Berlin.

\bibitem{ARSY2}
M. Alfonseca, D. Ryabogin, A. Stancu, V. Yaskin, ``\textit{Characterizations of balls and ellipsoids by infinitesimal homothetic conditions}'', Acta Math. Sci. Ser. B (Engl. Ed.) \textbf{45},  no. 1, 280--290 (2025).

\bibitem{BMO}
J. Bracho, L. Montejano, D. Oliveros, ``\textit{Carousels, Zindler curves and the floating body problem}", Period. Math. Hungar \textbf{49} (2004), 9–23.

%\bibitem{BP}
%H. Busemann, C.M. Petty, ``\textit{Problems on convex bodies}", Math. Scand., \textbf{4} (1956), 88--94.

\bibitem{Da}
A. Davidov, ``\textit{Theory of bodies floating in liquids in equilibrium}", Master Thesis, Moscow State University, 1848, found in Life and works of A. Yu. Davidov, (in Russian), by N.E. Zhukovsky,
P.A. Nekrasov and P.M. Pokrovsky, Math. Sb. \textbf{15} (1890), 1--57.

\bibitem{Du}
C. Dupin, ``\textit{Applications de G\'eom\'etrie et de M\'echanique \'a la marine, aux ponts et chaussées}", Paris (1822).

\bibitem{Hi}
M.W. Hirsch, ``\textit{Differential Topology}", Springer-Verlag (1976).

\bibitem{HSW}
H. Huang, B. Slomka, E. Werner, ``\textit{Ulam floating bodies}", J. of London Math. Soc. \textbf{100} (2019), 425--446.

\bibitem{KP}
S. Krantz, H. R. Parks, ``\textit{A Primer of Real Analytic Functions}", Birkhäuser (2002).

\bibitem{MM} Y. Martinez-Maure, ``\textit{Hedgehog Theory}'', hal-04404176v1f, available at \href{https://hal.science/hal-04404176v1/document}{https://hal.science/hal-04404176v1/document}.

\bibitem{MR}
M. Meyer, S. Reisner, ``\textit{A geometric property of the boundary of symmetric convex bodies and convexity of flotation surfaces}", Geom. Dedicata \textbf{37} (1991), 327–337.

\bibitem{MR2018}
S. Myroshnychenko, D. Ryabogin, “\textit{On polytopes with congruent projections or sections}", Advances in Mathematics, 2018 (\textbf{325}) 482--504.

\bibitem{Nu}
R. Nurnberg, ``\textit{Calculating the area and centroid of a polygon in 2d}'', Imperial College London (2013).

\bibitem{O}
S. P. Olovjanischnikoff, ``\textit{Ueber eine kennzeichnende Eigenschaft des Ellipsoides}", Leningrad State Univ. Ann. (Uchen. Zap.), \textbf{83} (1941), 113--128.

\bibitem{R22}
D. Ryabogin, ``\textit{A negative answer to Ulam’s problem 19 from the Scottish Book}", Ann. of Math. (2) \textbf{195} (2022), 1111--1150.

\bibitem{R23}
D. Ryabogin, ``\textit{On bodies floating in equilibrium in every orientation}", Geom. Dedicata (4) \textbf{217} (2023), Paper No. 70, 17 pp.


\bibitem{S}
R. Schneider, ``\textit{Convex Bodies: The Brunn–Minkowski Theory}", Encyclopedia of Mathematics and its Applications, 2nd ed. 2014.

\bibitem{SW}
C. Sch\"utt, E. Werner, ``\textit{The convex floating body}", Math. Scand. \textbf{66} (1990), 275--290.

\bibitem{VP}
C.-J. De la Vallée-Poussin, ``\textit{Leçons de Mécanique Analytique}", vol II, Paris, Gautier Villars Editeur, Louvain (1925).

\bibitem{W} 
E. Werner, ``\textit{Floating Bodies and Illumination Bodies}", Integral Geometry and Convexity, pp. 129--140 (2006).

\bibitem{Z}
N.E. Zhukovskii, ``\textit{Classical Mechanics}", Moscow (1936), in Russian.
\end{thebibliography}

\end{document}
